\begin{itemframe}{مقدمه}
\itm
بسیاری از مسائل به شکل ‌بهینه‌سازی یک تابع هدف تحت تعدادی قید مطرح می‌شوند.
\itm
اگر بتوانید تابع هدف را به صورت یک تابع خطی از متغیرها و قیود را به صورت معادلات یا نامعادلات خطی بر این متغیرها بیان کنید، آنگاه با یک «مسئلهٔ برنامه‌ریزی خطی»
\fn{linear programming problem}
روبرو هستید.
\end{itemframe}


\begin{itemframe}{مقدمه}
\itm
برای مثال بیایید یکی از کابردهای برنامه ریزی خطی را بررسی کنیم؛
فرض کنید شما یک سیاستمدار هستید که می‌خواهید در یک انتخابات پیروز شوید.
\itm
ناحیهٔ انتخاباتی شما شامل سه نوع منطقه است: شهری، حومه‌ای، و روستایی.
این مناطق، به ترتیب شامل 100، 200، و 50 هزار رأی‌دهندهٔ هستند.
هرچند همهٔ رأی‌دهندگان به پای صندوق رأی نمی‌روند، شما می‌خواهید دست‌کم نیمی از رأی‌دهندگان در هر منطقه به شما رأی دهند.
\end{itemframe}


\begin{itemframe}{مقدمه}
\itm
برخی موضوعات در مناطق خاصی در جلب رأی مؤثرتر اند. موضوعات موجود عبارتند از: آمادگی برای آخرالزمان زامبی‌ها، تجهیز کوسه‌ها با لیزر، ساخت بزرگراه برای خودروهای پرنده، حق رأی برای دلفین‌ها.
\itm
در جدول زیر نشان داده‌ شده که با صرف هر هزار تومان تبلیغات برای هر موضوع، چه تعداد رأی از هر قشر جمعیتی به‌دست می‌آورید یا از دست می‌دهید. مقادیر منفی بیانگر از دست دادن رأی هستند.
\pic{1}
\end{itemframe}


\begin{itemframe}{مقدمه}
\itm
حال مسئله این است: یافتن حداقل مقدار پول لازم برای جذب 50 هزار رأی شهری، 100 هزار رأی حومه‌ای، و 25 هزار رأی روستایی.
\itm
بیایید این مسئله را مدل سازی ریاضی کنیم. در گام نخست باید تصمیم‌هایی که شما باید در طول حل مسئله بگیرید بیابید و آنها را در غالب متغییرها معرفی کنید.
از آنجا که در این مسئله در مورد چهار مقدار باید تصمیم گرفته شود، چهار
«متغیر تصمیم»
\fn{decision variables}
 معرفی می‌کنیم:
$$
x_1, x_2, x_3, x_4
$$
این متغییرها به ترتیب تعداد هزاران تومان صرف‌شده برای موضوعات اول تا چهارم هستند.
\end{itemframe}


\begin{itemframe}{مقدمه}
\itm
حال بیایید قیود
\fn{constraints}
را مدل کنیم. به محدودیت‌های مقادیر متغیرها قید گفته می‌شود. شرط کسب دست‌کم 50,000 رأی شهری به صورت زیر بیان می‌شود:
$$
-2x_1 + 8x_2 + 0x_3 + 10x_4 \ge 50
$$
\itm
به‌طور مشابه، قیود برای حومه و روستا چنین هستند:
$$
5x_1 + 2x_2 + 0x_3 + 0x_4 \ge 100
$$

$$
3x_1 - 5x_2 + 10x_3 - 2x_4 \ge 25
$$
\end{itemframe}


\begin{itemframe}{مقدمه}
\itm
هر مقداردهی به متغیرهای
$x_1$
تا
$x_4$
که نامعادلات بالا را ارضا کند، تعداد رأی مورد نظر را کسب می‌کند.
\itm
در نهایت، به تابع هدف فکر کنید؛ کمیتی که می‌خواهید آن را بهینه کنید. در اینجا می‌خواهیم هزینه تبلیغات کمینه شود، پس تابع هدف به این صورت است:
$$
x_1 + x_2 + x_3 + x_4
$$
\itm
همچنین تبلیغات با هزینه منفی وجود ندارد. بنابراین باید داشته باشیم:
$$
x_1 \ge 0,\quad x_2 \ge 0,\quad x_3 \ge 0,\quad x_4 \ge 0
$$
\end{itemframe}


\begin{itemframe}{مقدمه}
\itm
ترکیب نامعادلات قید با تابع هدف، منجر به چیزی می‌شود که آن را یک
«برنامهٔ خطی»
\fn{linear program}
 می‌نامند:

%todo: text in math interprets inversely
\begin{align*}
minimize \quad & x_1 + x_2 + x_3 + x_4 \\
subject \: to \quad \\
& -2x_1 + 8x_2 + 0x_3 + 10x_4 \ge 50 \\
& 5x_1 + 2x_2 + 0x_3 + 0x_4 \ge 100 \\
& 3x_1 - 5x_2 + 10x_3 - 2x_4 \ge 25 \\
&x_1, x_2, x_3, x_4 \ge 0&
\end{align*}
\end{itemframe}


